\documentclass[10pt]{article}
\usepackage[T1]{fontenc}
\usepackage[utf8]{inputenc}
\usepackage[french]{babel}
\usepackage{amsmath,amssymb,mathrsfs,mathpazo}
\usepackage{geometry}
\geometry{a4paper,margin=2cm}
\usepackage{graphicx,array,tabularx,tabularray,colortbl,multirow,multicol,booktabs,enumitem,pifont,ulem,fancybox,diagbox,cancel}
\usepackage{pstricks,pstricks-add,pst-plot,pst-tree,pst-node,pst-text,pst-func,pst-math,pst-bspline,pst-3dplot,pst-eucl,pst-solides3d}
\let\pscalculate\undefined
\usepackage{pst-bezier}
\usepackage{tikz,pgfplots,tkz-tab}
\usetikzlibrary{arrows,arrows.meta,patterns,calc,positioning,trees,decorations.markings,shapes.geometric,shapes.misc,shapes.arrows,intersections,decorations.pathmorphing,decorations.pathreplacing,backgrounds,shadows,babel}
\pgfplotsset{compat=1.18}
\usepgfplotslibrary{fillbetween}
\usepackage[np]{numprint}
\usepackage{listings,makecell,pas-tableur,twoopt}
\let\barre\undefined
\usepackage[pstricks]{tablvar}
\definecolor{darkgreen}{RGB}{0,100,0}
\providecommand{\red}{\color{red}}
\providecommand{\blue}{\color{blue}}
\providecommand{\green}{\color{green}}
\providecommand{\black}{\color{black}}
\providecommand{\white}{\color{white}}
\providecommand{\gray}{\color{gray}}
\providecommand{\cyan}{\color{cyan}}
\providecommand{\magenta}{\color{magenta}}
\providecommand{\yellow}{\color{yellow}}
\providecommand{\violet}{\color{violet}}
\providecommand{\orange}{\color{orange}}
\providecommand{\euro}{\texteuro}
\providecommand{\EUR}{\texteuro}
\providecommand{\ieme}{\textsuperscript{e}}
\providecommand{\degres}{\ensuremath{^\circ}}
\providecommand{\degree}{\ensuremath{^\circ}}
\providecommand{\vv}[1]{\overrightarrow{#1}}
\providecommand{\Coord}[1]{#1}
\providecommand{\decosix}{}
\providecommand{\starredbullet}{\ensuremath{\bullet}}
\pagestyle{empty}
\setlength{\parindent}{0pt}
\usepackage{adjustbox}\DeclareUnicodeCharacter{2236}{:}\DeclareUnicodeCharacter{2190}{\ensuremath{\leftarrow}}
\providecommand{\R}{\mathbb{R}}
\providecommand{\N}{\mathbb{N}}
\providecommand{\Z}{\mathbb{Z}}
\providecommand{\Q}{\mathbb{Q}}
\providecommand{\D}{\mathbb{D}}
\providecommand{\C}{\mathbb{C}}
\providecommand{\e}{{\mathrm{e}}}
\providecommand{\vect}[1]{\overrightarrow{#1}}
\providecommand{\vectt}[1]{\overrightarrow{\mathrm{#1}}}
\providecommand{\vvt}[1]{\overrightarrow{\mathrm{#1}}}
\providecommand{\np}[2][]{{#2~#1}}
\providecommand{\vv}[1]{{\overrightarrow{#1}}}
\renewcommand{\d}{{\mathrm{d}}}
\providecommand{\degres}{{^\circ}}
\providecommand{\degre}{{^\circ}}
\providecommand{\blue}{\color{blue}}
\providecommand{\red}{\color{red}}
\providecommand{\enskip}{\quad}
\providecommand{\llbracket}{\lbrack\!\lbrack}
\providecommand{\rrbracket}{\rbrack\!\rbrack}
\providecommand{\euro}{{\text{€}}}
\providecommand{\no}{{\text{n°}}}
\providecommand{\interval}[2]{{[#1;#2]}}
\providecommand{\intervFO}[2]{{[#1;#2[}}
\providecommand{\dotfill}{\ldots}

\providecommand{\ds}{\displaystyle}
\providecommand{\textding}[1]{#1}
\providecommand{\Oij}{\(\mathrm{O};\vec{\imath},\vec{\jmath}\)}
\providecommand{\Oijk}{\(\mathrm{O};\vec{\imath},\vec{\jmath},\vec{k}\)}
\providecommand{\Ouv}{\(\mathrm{O};\vec{u},\vec{v}\)}
\providecommand{\intervFF}[2]{\ensuremath{[#1;#2]}}
\renewcommand{\euro}{\ifmmode\text{\texteuro}\else\texteuro\fi}
\DeclareUnicodeCharacter{00B0}{\ifmmode^{\circ}\else\textdegree\fi}
\DeclareUnicodeCharacter{20AC}{\ifmmode\text{\texteuro}\else\texteuro\fi}
% Énoncés : Éducation nationale. Transcription et corrigés : APMEP, Denis Vergès et contributeurs. Découpage et classement : CoopMaths, Sébastien Lozano, Éric Elter et contributeurs. Adaptation de la présentation : Maths Explorer.
% CC BY-SA : https://coopmaths.fr/www/about/
% https://coopmaths.fr/alea/static/bac/2026/tex/bac_2026_06_sujet2_polynesie_4.tex
\begin{document}
\noindent\textit{APMEP / CoopMaths. CC BY-SA. Présentation adaptée par Maths Explorer.}\par\medskip
Pour chacune des affirmations suivantes, indiquer si elle est vraie ou fausse.

\emph{Justifier chaque réponse. Une réponse non justifiée ne rapporte aucun point.}

\medskip

\begin{enumerate}
\item On considère la fonction $f$ définie pour tout $x \in \mathbb{R}$ par $f(x) = x \e^{(x^2)}$.

\textbf{Affirmation 1} : $f$ est concave sur $]-\infty~;~0]$.

\item Un refuge spécialisé accueille 24 rapaces numérotés de 1 à 24. Le refuge travaille avec trois vétérinaires. Afin d'administrer un soin à tous les oiseaux, le directeur veut donc former trois groupes, numérotés de 1 à 3, constitués chacun de 8 rapaces.

\textbf{Affirmation 2} : Le nombre de façons de répartir les rapaces dans les groupes 1 à 3 est :

\[\frac{24!}{(8!)^3}.\]

\item Soit $f$ une fonction définie et dérivable sur $\mathbb{R}$, solution de l'équation différentielle :

\[y' = -y + 1 \text{ sur } \mathbb{R}.\]

\textbf{Affirmation 3} : La fonction $g$ définie pour tout $x \in \mathbb{R}$ par $g(x) = f(x) - 1$ vérifie

\[g' = -g.\]

\item La suite $(u_n)$ est définie pour tout entier naturel $n$ non nul par :

\[u_n = \frac{\sin(n)}{n}.\]

\textbf{Affirmation 4} : La suite $(u_n)$ n'admet pas de limite.

\item On a tracé ci-dessous la parabole d'équation $y = -x^2 + 4x + 1$ et la droite d'équation $y = x + 1$.

\begin{center}


















\psset{unit=1cm,arrowsize=2pt 3}
\begin{pspicture*}(-1.1,-1.5)(5.5,5.5)
\psaxes[linewidth=1.25pt,labelFontSize=\scriptstyle]{->}(0,0)(-1,-1.5)(5.5,5.5)
\psplot[plotpoints=2000,linewidth=1.25pt,linecolor=red]{-0.5}{4.5}{4 x mul 1 add x dup mul sub}
\psplot[plotpoints=500,linewidth=1.25pt,linecolor=blue]{-1}{5}{x  1 add}
\pscustom[fillstyle=solid,fillcolor=lightgray]{\psplot[plotpoints=2000,linewidth=1.25pt,linecolor=red]{0}{3}{4 x mul 1 add x dup mul sub}\psline(3,4)(0,1)}
\psgrid[gridlabels=0pt,subgriddiv=2,gridwidth=0.2pt,subgridwidth=0.1pt]
\end{pspicture*}

\end{center}

\textbf{Affirmation 5} : L'aire du domaine grisé est égale à $6{,}5$ unités d'aire.
\end{enumerate}

\end{document}
